\section{Conclusions\label{sec:conclusions}}
This study built experience profiles based on the activity of 71,467 GitHub developers linked to the AIDev dataset, and examined how these profiles relate to the intensity of using AI coding agents and to pull-request acceptance, alongside a separate comparison of agent and human contribution acceptance. None of the three research questions shows a simple, monotonic relationship with activity level. Agent-use intensity is clearly non-monotonic. The Mid profile reports far more Agentic-PRs than both Junior and the highest-activity profile, Senior, and this is the only result with a large effect size (RQ1). The acceptance rate is highest for Mid and lowest for Senior, with Junior in between (RQ2). Agentic-PRs are accepted slightly more often than the independent baseline human population from a comparably popular subset of repositories (RQ3). In short, higher combined activity on GitHub is not a reliable sign of better or more frequent outcomes when using a coding agent: the highest-activity profile - Senior - has neither the most Agentic-PRs nor the highest acceptance rate among the three profiles.

We ran the sensitivity checks flagged in Section 3 with respect to an alternative $k$, repeated random initializations, kept outliers, exclusion of AIDev-derived indicators, and the complete-case restriction (Section 5). Their results were mixed. The clustering is very stable with respect to different random seeds (ARI = 0.98-0.99) and fairly stable with respect to the choice of $k$, but clearly less stable with respect to whether outliers were removed (ARI 0.64), and least stable with respect to whether the three agent-use indicators were included (ARI 0.32), which specifically weakens the RQ1 effect as well. This last result should temper how we read RQ1's large effect size: part of this effect likely reflects the fact that the clustering was partly built from the same kind of agent-use signal that RQ1 measures as an outcome. RQ1 and RQ2 proved robust to the complete-case restriction - their group means and pattern of significance were nearly unchanged after excluding the 1755 developers with estimated data. It is worth noting that we only re-ran the actual RQ1/RQ2 outcome statistics for two of the four checks - the complete-case restriction and the indicator exclusion. The alternative-$k$ and kept-outliers checks assessed the stability of the clustering itself, not whether the RQ1/RQ2 conclusions hold under them. The direction of RQ1 is consistent across all the checks performed, but its size is the most sensitive to design choices.

The bootstrap confidence intervals around each reported effect size (Section 4) are consistent with the picture above: the intervals for the pairwise comparisons in RQ1 and RQ2 are narrow relative to their point estimates, while the interval for RQ3 is comparatively wider, reflecting the much smaller and unevenly sized human-authored comparison group. Future work should extend RQ3 as more data on agent use is collected for a larger, more representative set of human-authored baseline repositories, rather than only the currently available high-star sample, which would allow the comparison with the human baseline to be assessed with more balanced group sizes than the current analysis allows. More broadly, as coding agents and the norms for reviewing their output keep changing, these relationships reflect one, current stage in the adoption of these tools, rather than a stable feature of human-agent collaboration on GitHub.
